\chapter{Quick reference}\label{app:quickref}

\section{Commands}

\begin{shell}
source /usr/lib/openfoam/openfoam2412/etc/bashrc

# geometry and one mesh
python3 mesh/makeRocket.py mesh/stl/Aconcagua_biconvex117_roll45.stl --finRoll 45
mesh/Allmesh ~/meshes/h1 --maxCellSize 1.6 --nLayers 10 --thicknessRatio 1.2
./post/firstLayer.py ~/meshes/h1

# one case
./newCase.sh ~/runs/test --mesh ~/meshes/h1 --solver pimple --Minf 0.8 --np 8
cd ~/runs/test && ./Allrun 8

# sweep
./runMesh.sh meshes.txt ~/meshes
nohup ./runCase.sh sweepSubsonic.nemepad.txt > ~/log.sweepSubsonic.nemepad 2>&1 &

# results
./post/monitor.sh ~/runs/test
./post/sweepTable.sh ~/runs > sweep.csv
python3 flightData/cdFlight.py

# check the binding and the measurement of a running case
for p in $(pgrep -x rhoPimpleFoam); do taskset -cp $p; done
tail -3 ~/runs/<case>/log.machine
\end{shell}

\section{Run names}

\texttt{m<M>-a<alpha>-h<n>[-pimple]}: \texttt{m08} is $M = 0.8$, \texttt{a04} is
$\alpha = 4^\circ$, \texttt{h1} is the mesh. The suffix \texttt{-pimple} marks the
\cmd{rhoPimpleFoam} start run of a supersonic chain.

\section{Key values}

\begin{tabularx}{\linewidth}{@{}P{0.42\linewidth} L@{}}
\toprule
Item & Value \\
\midrule
Reference area, full rocket & $A_\text{ref} = \pi\,(0.0755)^2 = \unit{0.0179081}{m\textsuperscript{2}}$ \\
Reference length & $l_\text{ref} = d = \unit{0.151}{m}$ \\
Moment reference & Nose tip, $(0, 0, 0)$ \\
Rocket length & \unit{2.955}{m} \\
First cell on the body & $t_1 \approx \unit{0.10}{mm}$, $y \approx \unit{0.05}{mm}$ \\
Stop criterion & $|\overline{C_D} - C_D| < 0.007$ over 1000 iterations \\
Memory & approximately \unit{3}{GB} for each million cells \\
Midland atmosphere & $p_0 = \unit{90\,565}{Pa}$ at \unit{896}{m} MSL, $T_0 = \unit{305.0}{K}$, \unit{7.2}{K/km} \\
\bottomrule
\end{tabularx}
